\SourceLeaf{311}%
% A544-P0311-D01
\[
(-1)^{\frac{p-1}4-\mu}2^{\frac{p-1}4}\equiv AB\equiv f\pmod p
\]
und
\[
2^{\frac{p-1}4}\equiv(-1)^{\mu-1}f\equiv f^{2\mu-1}\equiv f^{m+\frac{p-5}4}\pmod p
\]
da endlich in diesem Fall $b\equiv2\pmod4$, und folglich
\[
\frac{p-5}4=(a+1)\frac{a-1}4+\frac{bb-4}4\equiv2\frac{a-1}4\equiv\frac{a-1}2\pmod4
\]
ist, so erhält man wieder die zu beweisende Congruenz
\[
2^{\frac{p-1}4}\equiv f^{m+\frac{a-1}2}\pmod p
\]
\BookRule{35mm}

% A544-P0311-D02
\pdfbookmark[2]{Zu IX.}{nexus-remarks-ix}
\begin{center}\large Zu IX.\end{center}

% A544-P0311-D03
Es sei $p$ eine positive ungerade durch kein Quadrat theilbare Zahl, und
\[
S_r=\Sigma\left(\frac{s_r}p\right)
\]
wo $s_r$ alle relativen Primzahlen zu $p$ durchlaufen muss, welche zwischen $(r-1)\frac p8$ und $r\frac p8$ liegen; bezeichnet man die Anzahlen der nicht eigentlich äquivalenten formae proprie primitivae für die Determinanten $-p$ und $-2p$ resp. mit $C_1$ und $C_2$, so ist (vergl. \textsc{Dirichlet} Recherches sur diverses applications etc. §. 11 in \textsc{Crelle}'s Journal XXI)
\[
C_1=2(S_1+S_2),\quad C_2=2(S_1-S_4)
\]
oder
\[
C_1=S_1+S_2+S_3+S_4,\quad C_2=2(S_2+S_3)
\]
je nachdem $p\equiv1$ oder $\equiv3\pmod4$ ist. Bedenkt man ferner, dass die Zahlen $s_1$ und $s_2$ im Complex mit den Zahlen $2s_1$ und $p-2s_4$, und ebenso die Zahlen $s_3$ und $s_4$ im Complex mit den Zahlen $2s_2$ und $p-2s_3$ übereinstimmen, und dass im Falle $p\equiv1\pmod4$ die Summe $S_1+S_2+S_3+S_4=0$ ist, so ergeben sich in beiden Fällen noch zwei neue Relationen zwischen den vier Summen $S_1,S_2,S_3,S_4$, so dass jede derselben durch $C_1$ und $C_2$ ausgedrückt werden kann. Man erhält auf diese Weise, wenn $p\equiv1\pmod4$ ist,
\[
\begin{aligned}
S_1=S_8&=\tfrac14\left(\frac2p\right)C_1+\tfrac14C_2\\
S_2=S_7&=\tfrac14\left(2-\left(\frac2p\right)\right)C_1-\tfrac14C_2\\
S_3=S_6&=-\tfrac14\left(2+\left(\frac2p\right)\right)C_1+\tfrac14C_2\\
S_4=S_5&=\tfrac14\left(\frac2p\right)C_1-\tfrac14C_2
\end{aligned}
\]
\SourceLeaf{312}%
% A544-P0312-D01
und, wenn $p\equiv3\pmod4$ ist
\[
\begin{aligned}
S_1=-S_8&=\tfrac14\left(3+\left(\frac2p\right)\right)C_1-\tfrac14C_2\\
S_2=-S_7&=-\tfrac14\left(1-\left(\frac2p\right)\right)C_1+\tfrac14C_2\\
S_3=-S_6&=\tfrac14\left(1-\left(\frac2p\right)\right)C_1+\tfrac14C_2\\
S_4=-S_5&=\tfrac14\left(1-\left(\frac2p\right)\right)C_1-\tfrac14C_2
\end{aligned}
\]

% A544-P0312-D02
Ist $p$ eine Primzahl, so findet man hieraus unmittelbar die im Text angegebenen Formeln für die Anzahl der quadratischen Reste, welche in den einzelnen Octanten enthalten sind.
\BookRule{35mm}

% A544-P0312-D03
\pdfbookmark[2]{Zu X.}{nexus-remarks-x}
\begin{center}\large Zu X.\end{center}

% A544-P0312-D04
Es sei $p$ eine positive und durch kein Quadrat theilbare Zahl von der Form $6n\pm1$, und
\[
S_r=\Sigma\left(\frac{s_r}p\right)
\]
wo $s_r$ alle relativen Primzahlen zu $p$ durchlaufen muss, welche zwischen $(r-1)\frac p{12}$ und $r\frac p{12}$ liegen; bezeichnet man die Anzahlen der nicht eigentlich äquivalenten formae proprie primitivae für die Determinanten $-p$ und $-3p$ mit $C_1$, $C_3$, so findet man leicht (vergl. \textsc{Dirichlet} Recherches etc. §. 11 oder die Note zu VI und VII)
\[
C_1=2(S_1+S_2+S_3),\quad C_3=2(S_1+S_2-S_5-S_6)
\]
oder
\[
C_1=S_1+S_2+S_3+S_4+S_5+S_6,\quad C_3=2(S_2+S_3+S_4+S_5)
\]
je nachdem $p\equiv1$ oder $\equiv3\pmod4$ ist. Berücksichtigt man ferner, dass
\begin{center}\small
\begin{tabular}{@{}ccccccl@{}}
die&Zahlen&$s_1$ und $s_2$&mit&den&Zahlen&$2s_1$ und $p-2s_6$\\
„&„&$s_3$ und $s_4$&„&„&„&$2s_2$ und $p-2s_5$\\
„&„&$s_5$ und $s_6$&„&„&„&$2s_3$ und $p-2s_4$
\end{tabular}
\end{center}
und ebenso
\begin{center}\small
\begin{tabular}{@{}ccccccl@{}}
die&Zahlen&$s_1,s_2,s_3$&mit&den&Zahlen&$3s_1,\ 3s_5-p,\ p-3s_4$\\
„&„&$s_4,s_5,s_6$&„&„&„&$3s_2,\ 3s_6-p,\ p-3s_3$
\end{tabular}
\end{center}
übereinstimmen, und dass im Falle $p\equiv1\pmod4$ die Summe $S_1+S_2+S_3+S_4+S_5+S_6=0$ ist, so erhält man ausser den beiden obigen noch vier neue Relationen zwischen den sechs Summen $S_1,S_2\ldots S_6$, so dass dieselben sämmtlich aus $C_1$ und $C_3$ bestimmt werden können. Man erhält auf diese Weise, wenn $p\equiv1\pmod4$ ist,
\SourceLeaf{313}%
% A544-P0313-D01
\[
\begin{aligned}
S_1=S_{12}&=\tfrac14\left(1+\left(\frac3p\right)\right)C_1+\tfrac1{12}\left(1+\left(\frac2p\right)\right)C_3\\
S_2=S_{11}&=-\tfrac14\left(1+\left(\frac3p\right)\right)C_1+\tfrac1{12}\left(1+\left(\frac2p\right)\right)C_3\\
S_3=S_{10}&=\tfrac12C_1-\tfrac16\left(1+\left(\frac2p\right)\right)C_3\\
S_4=S_9&=-\tfrac12C_1-\tfrac16\left(2-\left(\frac2p\right)\right)C_3\\
S_5=S_8&=\tfrac14\left(1+\left(\frac3p\right)\right)C_1-\tfrac1{12}\left(5-\left(\frac2p\right)\right)C_3\\
S_6=S_7&=-\tfrac14\left(1+\left(\frac3p\right)\right)C_1+\tfrac1{12}\left(1+\left(\frac2p\right)\right)C_3
\end{aligned}
\]
und, wenn $p\equiv3\pmod4$ ist,
\[
\begin{aligned}
S_1=-S_{12}&=\tfrac1{12}\left(9+3\left(\frac2p\right)-\left(\frac3p\right)+\left(\frac6p\right)\right)C_1-\tfrac14C_3\\
S_2=-S_{11}&=\tfrac14\left(-1+\left(\frac2p\right)+\left(\frac3p\right)-\left(\frac6p\right)\right)C_1+\tfrac14C_3\\
S_3=-S_{10}&=\tfrac16\left(-\left(\frac3p\right)+\left(\frac6p\right)\right)C_1\\
S_4=-S_9&=\tfrac16\left(3-2\left(\frac3p\right)-\left(\frac6p\right)\right)C_1\\
S_5=-S_8&=\tfrac14\left(-1-\left(\frac2p\right)+\left(\frac3p\right)+\left(\frac6p\right)\right)C_1+\tfrac14C_3\\
S_6=-S_7&=\tfrac1{12}\left(3-3\left(\frac2p\right)+\left(\frac3p\right)-\left(\frac6p\right)\right)C_1-\tfrac14C_3
\end{aligned}
\]

% A544-P0313-D02
Ist $p$ eine Primzahl, so findet man aus dem ersten System die im Text angegebenen Formeln; für die andern Fälle erhält man ähnliche Formeln aus dem zweiten System.
\begin{flushright}\textsc{R. Dedekind.}\end{flushright}
\BookRule{55mm}

\clearpage
\renewcommand{\BookRunningTitle}{Geometrische Seite der ternären Formen}
\SourceLeaf{315}%
% A544-P0315-D01
\pdfbookmark[0]{Geometrische Seite der ternären Formen}{ternary-geometry}
\begin{center}\Large\bfseries GEOMETRISCHE SEITE DER TERNÄREN FORMEN.\end{center}
\BookRule{30mm}

% A544-P0315-D02
Ein Punkt im Raume $(0)$ sei als Anfangspunkt angenommen. Der Uebergang von da zu drei andern Punkten $P,P',P''$, die mit jenem nicht in einer Ebene liegen, sei resp. $t,t',t''$; wo, so oft keine Verwechslung möglich ist, die Punkte $P,P',P''$ selbst durch $(t)$. $(t')$, $(t'')$ bezeichnet werden mögen.

% A544-P0315-D03
Es sei ferner allgemein $(t,t')$ das Product der Länge der beiden Linien $t,t'$ in den Cosinus ihrer Neigung etc.

% A544-P0315-D04
Man hat allgemein $(\alpha t+\alpha't'+\alpha''t''+\ldots\quad\beta u+\beta'u'+\beta''u''+\mathinner{.\,.})$ wenn man die Multiplication
\[
(\alpha t+\alpha't'+\alpha''t''+\mathinner{.\,.})\times(\beta u+\beta'u'+\beta''u''+\mathinner{.\,.})
\]
ausführt und statt $tu,\allowbreak tu',\allowbreak tu'',\allowbreak t'u',\allowbreak t'u''\mathinner{.\,.}$ u. s. w. $(t,u),\allowbreak(t,u'),\allowbreak(t,u''),\allowbreak(t',u'),\allowbreak(t',u'')$ u. s. w. schreibt.

% A544-P0315-D05
Jeder Punkt im Raume wird durch ein Trinomium
\[
(xt+x't'+x''t'')
\]
dargestellt werden können.

% A544-P0315-D06
Für alle Punkte, die in einer bestimmten Ebene liegen, wird dann eine Gleichung
\[
\lambda x+\lambda'x'+\lambda''x''=L
\]
